\section{Rをつかった項目反応理論}
\subsection{授業内容}
\subsubsection{科目の中でのこのコマの位置づけ}
項目反応理論を実践的に理解する演習パートである。

カテゴリカルな因子分析と数学的に同等ではあるが，より項目の特徴を広く表現できる項目反応理論の利用が，今後より重要なものになってくるだろう。

ここではまずテスト理論の根本に立ち返り，二値単因子のデータを使って1PL,2PLモデルの分析を行う。分析結果は数値で見ることも重要であるが，ICCやIIC，TICなどを使って可視化するとより理解が深まるだろう。
多段階の反応についても，同様にGRMを実行し，閾値や識別力，IRCCCやIIC，TICが描画できることを確認する。特にIRCCCによる反応段階の読み取り方には注意する。
最後に多段階で多因子の場合，項目反応理論の文脈から言えば多次元IRTになり専用のパッケージが必要になることを紹介しつつ，カテゴリカル因子分析でも同様のことができることを確認する。
\subsubsection{コマ主題細目}
\begin{description}
	\item[項目反応理論の実際]項目反応理論はテスト理論がその出自に当たるので，まずは二値データで単因子が想定できるような例を元に分析を行う。分析には\texttt{irtoys}パッケージや\texttt{ltm}パッケージを用いて，1PLモデル，2PLモデルの演習を行う。項目母数の値と意味が，具体的な設問に照らし合わせて考えることで，より実感をもって理解できるようになると思われる。特に，ICCやIIC,TICなど可視化することでその意味が理解しやすくなるだろう。3PLモデルなどさらに拡張したモデルも利用可能である。

	\item[段階反応モデルの実際] 続いて単因子，段階反応モデルの実践を行う。因子構造として，前回の授業で扱った多因子の内，ある因子に限定して分析を行うこととする。段階反応モデル(GRM)の出力結果を数値だけでなく可視化することで，項目の特徴がどのように表現されているかを考える。特に反応段階の山が潰れているようなケースは，適切な反応段階でなかったことを意味するので，数値の置き換えなど元データを修正しつつ分析し直すことを考える。これらを通じて，適切な反応段階による調査法が必要であることを理解する。

	\item[カテゴリカル因子分析との対応] 多段階，多因子の場合は\texttt{psych}パッケージの\texttt{fa}関数にオプションを追加することでできる，カテゴリカル因子分析と同じである。出力結果について，これまでの相関係数を用いているものとの違いを確認する。また数値をどのように変換すれば対応するのかを見ることで，数学的に等価であることを確認しておく。IRTの側面から多因子に拡張した，多次元IRTについても，\texttt{mirt}パッケージを利用すれば実行できる。この解析には計算時間がかかるが，完全情報最尤推定の結果が得られることは利点である。
\end{description}
\subsubsection{キーワード}
\begin{itemize}
	\item 1PLモデル，2PLモデル，3PLモデル
	\item 段階反応モデル
	\item カテゴリカル因子分析
	\item \texttt{irtoys,ltm,psych,mirt}パッケージ
\end{itemize}
\subsection{授業情報}
\paragraph{コマの展開方法}
Rを使った演習
\subsubsection{予習・復習課題}
\paragraph{予習} \texttt{irtoys,ltm}パッケージを事前にインストールして環境を整え，データファイルの読み込みなどR/RStudioの基本的な使い方を確認しておこう。特にRの\texttt{data.frame}型に含まれる変数が，\texttt{numeric}なのか\texttt{factor}なのかによって挙動がかわることがある。変数の型についても再確認しておこう。
\paragraph{復習}本講で習ったパッケージを使って，具体的なデータを因子分析，IRT，カテゴリカルIRTなどいくつかのモデルで分析し，それぞれの違いを確認しておこう。